\hypertarget{chapter-05-page-103}{}
\subsection{Calderón--Zygmund 定理}
\label{sec:ch5-calderon-zygmund-theorem}

上一章用 Hilbert 变换研究了奇异积分. 本章考察这样的奇异积分: 它们的核具有 Hilbert 变换核的基本性质. 这使我们可以推广定理~\ref{thm:ch3-riesz-kolmogorov}, 得到如下结果.

\begin{Theorem}[Calderón--Zygmund]
\label{thm:ch5-calderon-zygmund}
设 $K$ 是 $\mathbb R^n$ 上的缓增分布, 它在 $\mathbb R^n\setminus\{0\}$ 上与局部可积函数一致, 并且
\begin{equation}
\label{eq:ch5-bounded-fourier-kernel}
 |\widehat K(\xi)|\le A,
\end{equation}
\begin{equation}
\label{eq:ch5-hormander-condition}
 \int_{|x|>2|y|}|K(x-y)-K(x)|\,dx\le B,
 \qquad y\in\mathbb R^n.
\end{equation}
则对 $1<p<\infty$,
\[
 \|K*f\|_p\le C_p\|f\|_p,
\]
并且
\[
 |\{x\in\mathbb R^n:|K*f(x)|>\lambda\}|
 \le\frac C\lambda\|f\|_1.
\]
\end{Theorem}

先对 $f\in\mathcal S$ 证明这些不等式, 再像处理 Hilbert 变换那样推广到任意 $f\in L^p$. 条件 \eqref{eq:ch5-hormander-condition} 通常称为 Hörmander 条件. 实际应用中, 它常由更强的梯度条件推出.

\hypertarget{chapter-05-page-104}{}
\begin{Proposition}
\label{prop:ch5-gradient-implies-hormander}
若对每个 $x\ne0$ 都有
\begin{equation}
\label{eq:ch5-gradient-condition}
 |\nabla K(x)|\le\frac C{|x|^{n+1}},
\end{equation}
则 Hörmander 条件 \eqref{eq:ch5-hormander-condition} 成立.
\end{Proposition}

这由中值定理直接推出, 细节留给读者.

\begin{Proof}[定理~\ref{thm:ch5-calderon-zygmund} 的证明]
这个证明本质上重复定理~\ref{thm:ch3-riesz-kolmogorov} 的证明, 因而省略部分细节. 令 $f\in\mathcal S$, $Tf=K*f$. 由 \eqref{eq:ch5-bounded-fourier-kernel}, $\|Tf\|_2\le A\|f\|_2$. 只需证明 $T$ 是弱型 $(1,1)$ 的: $1<p<2$ 的强型不等式由插值得到, $p>2$ 的情形则由对偶性得到, 因为伴随算子 $T^*$ 的核 $K^*(x)=\overline{K(-x)}$ 也满足 \eqref{eq:ch5-bounded-fourier-kernel} 与 \eqref{eq:ch5-hormander-condition}.

固定 $\lambda>0$, 对 $f$ 作高度为 $\lambda$ 的 Calderón--Zygmund 分解. 写 $f=g+b$, 其中 $g\in L^2$, $|g|\le2^n\lambda$, 而 $b=\sum_jb_j$, 每个 $b_j$ 支于两两不交的立方体 $Q_j$ 且积分为零. 与前述证明一样, 问题归结为
\begin{equation}
\label{eq:ch5-bad-part-outside-cubes}
 \int_{\mathbb R^n\setminus Q_j^*}|Tb_j(x)|\,dx
 \le C\int_{Q_j}|b_j(x)|\,dx,
\end{equation}
其中 $Q_j^*$ 与 $Q_j$ 同心, 边长为后者的 $2\sqrt n$ 倍. 记公共中心为 $c_j$. 对 $x\notin Q_j^*$, 利用 $\int b_j=0$ 得
\[
 Tb_j(x)=\int_{Q_j}[K(x-y)-K(x-c_j)]b_j(y)\,dy.
\]
由于 $\mathbb R^n\setminus Q_j^*\subset\{x:|x-c_j|>2|y-c_j|\}$, 对积分次序作交换并应用 \eqref{eq:ch5-hormander-condition}, 即得 \eqref{eq:ch5-bad-part-outside-cubes}.
\end{Proof}

现在考察上一章中的 $-n$ 次齐次核 $K(x)=\Omega(x')/|x|^n$. 要使 Hörmander 条件成立, 对 $\Omega$ 应作什么假设? 命题~\ref{prop:ch5-gradient-implies-hormander} 表明 $\Omega\in C^1(S^{n-1})$ 足够, 但还可以采用弱得多的条件. 定义
\[
 \omega_\Omega(t)=\sup\{|\Omega(u_1)-\Omega(u_2)|:
 |u_1-u_2|<t,\ u_1,u_2\in S^{n-1}\}.
\]

\hypertarget{chapter-05-page-105}{}
\begin{Proposition}
\label{prop:ch5-dini-implies-hormander}
若 $\Omega$ 满足
\begin{equation}
\label{eq:ch5-dini-condition}
 \int_0^1\omega_\Omega(t)\frac{dt}{t}<\infty,
\end{equation}
则核 $\Omega(x')/|x|^n$ 满足 \eqref{eq:ch5-hormander-condition}.
\end{Proposition}

条件 \eqref{eq:ch5-dini-condition} 称为 Dini 型条件.

\begin{Proof}
写
\[
 |K(x-y)-K(x)|
 \le\frac{|\Omega((x-y)')-\Omega(x')|}{|x-y|^n}
 +|\Omega(x')|\left|\frac1{|x-y|^n}-\frac1{|x|^n}\right|.
\]
条件 \eqref{eq:ch5-dini-condition} 蕴含 $\Omega$ 有界. 在 $|x|>2|y|$ 上, 第二项由 $C|y|/|x|^{n+1}$ 控制, 因而可积. 同时 $|(x-y)'-x'|\le4|y|/|x|$, 所以第一项在该区域的积分不超过
\[
 2^n|S^{n-1}|\int_0^2\omega_\Omega(t)\frac{dt}{t}<\infty.
\]
\end{Proof}

类似结果见第~\ref{subsec:ch5-general-dini}~节. 命题~\ref{prop:ch5-dini-implies-hormander} 与推论~\ref{cor:ch4-bounded-multiplier-kernel} 立即给出定理~\ref{thm:ch5-calderon-zygmund} 的如下推论.

\begin{Corollary}
\label{cor:ch5-dini-singular-integral}
若 $\Omega$ 定义在 $S^{n-1}$ 上, 积分为零并满足 \eqref{eq:ch5-dini-condition}, 则算子
\[
 Tf(x)=\operatorname{p.v.}\int_{\mathbb R^n}
 \frac{\Omega(y')}{|y|^n}f(x-y)\,dy
\]
是强型 $(p,p)$ 的, $1<p<\infty$, 并且是弱型 $(1,1)$ 的.
\end{Corollary}

\hypertarget{chapter-05-page-106}{}
由于 \eqref{eq:ch5-dini-condition} 蕴含 $\Omega$ 有界, 强型 $(p,p)$ 不等式也可由旋转方法得到. 新得到的是弱型 $(1,1)$ 不等式.
